\documentclass[aspectratio=169]{beamer}

% --- PAQUETES ---
\usepackage[utf8]{inputenc}
\usepackage[spanish]{babel}
\usepackage{graphicx}
\usepackage{booktabs}
\usepackage{multicol}
\usepackage{amsmath}
\usepackage{amssymb}
\usepackage{tikz}
\usetikzlibrary{positioning, arrows.meta, shapes, fit, calc}
\usepackage{pgfplots}
\pgfplotsset{compat=1.17}
\usepackage{hyperref}
\usepackage{url}

% --- TEMA Y ESTILO ---
\usetheme{Madrid}
\usecolortheme{default}

% Definición de colores (misma paleta que las presentaciones de Venegas)
\definecolor{DeepBlue}{RGB}{0,51,102}
\definecolor{AlertRed}{RGB}{204,0,0}
\definecolor{SafeGreen}{RGB}{0,100,0}
\definecolor{OrangeAccent}{RGB}{230,126,34}

% --- METADATA ---
\title[Redes Duales BID/CORFO]{Redes de Financiamiento Duales: \\ Aplicando el Modelo de Jalan et al. (2024) \\ al Caso BID/CORFO}
\subtitle{Validación empírica del modelo dual de Venegas (2025)}
\author{Eduardo Caballero \\ \small{Profesor Guía: Sebastián Cea}}
\institute{Universidad de Los Andes}
\date{\today}

\begin{document}

% ===================================================================
% PORTADA
% ===================================================================
\frame{\titlepage}

\begin{frame}{Contenidos}
\tableofcontents
\end{frame}

% ===================================================================
% SECCIÓN 1: MOTIVACIÓN
% ===================================================================
\section{Motivación}

\begin{frame}{El problema: financiamiento a MiPyMEs en Chile}
\textbf{Contexto:}
\begin{itemize}
    \item Las micro y pequeñas empresas (MiPyMEs) enfrentan restricciones de acceso al crédito bancario tradicional.
    \item Desde 2016, con apoyo de un préstamo del \textcolor{blue}{BID (US\$120 millones)}, CORFO opera el programa \textbf{Crédito Corfo Mipyme}: banca de segundo piso a través de intermediarios financieros no bancarios (IFNB).
    \item La MiPyME \textbf{no puede acceder directamente} a los fondos CORFO/BID: debe pasar por un intermediario (cooperativa, leasing, factoring).
\end{itemize}

\vspace{0.3cm}
\textbf{Pregunta:} ¿Cómo compite este canal público-regulado con el canal bancario tradicional por el financiamiento de la MiPyME, y qué le ocurre al bienestar agregado cuando el acceso directo está bloqueado?

\vspace{0.3cm}
\alert{Enfoque:} modelar el sistema como una \textbf{red financiera} en la que cada institución elige su cartera de contratos para maximizar utilidad, sujeta a restricciones regulatorias.
\end{frame}

\begin{frame}{Por qué un modelo de redes financieras}
\begin{itemize}
    \item La intermediación obligatoria (CORFO $\not\to$ MiPyME directo) es exactamente el tipo de \textcolor{blue}{restricción de red} que estudian Jalan, Chakrabarti \& Sarkar (2024).
    \item Venegas (2025) extiende ese modelo a \textbf{dos redes en competencia} (tradicional vs. alternativa) para estudiar la emergencia de infraestructuras de financiamiento no tradicionales.
    \item Este trabajo toma esa extensión dual y la \textbf{lleva a datos reales} del caso BID/CORFO, en vez de parámetros ilustrativos.
\end{itemize}

\vspace{0.4cm}
\begin{block}{Objetivo del trabajo de título}
Calibrar y validar el modelo dual de Venegas con datos públicos de CORFO/BID, la banca tradicional (CMF) y encuestas de empresas (INE), para obtener contratos, precios y reparto de bienestar consistentes con el caso chileno.
\end{block}
\end{frame}

% ===================================================================
% SECCIÓN 2: MODELO BASE (JALAN ET AL., 2024)
% ===================================================================
\section{Modelo Base (Jalan et al., 2024)}

\begin{frame}{Estructura de la red}
\textbf{Representación:} $n$ agentes conectados por una red ponderada $W \in \mathbb{R}^{n\times n}$.

\vspace{0.3cm}
\begin{itemize}
    \item \textbf{Tamaño de contrato:} $W_{ij}$ = monto acordado entre agentes $i$ y $j$.
    \begin{itemize}
        \item Simétrico: $W_{ij}=W_{ji}$ (requiere acuerdo mutuo).
    \end{itemize}
    \item \textbf{Precio del contrato:} $P_{ji}=-P_{ij}$ (suma cero, pagos entre partes).
\end{itemize}

\vspace{0.3cm}
Los contratos \textbf{emergen endógenamente} de la interacción de los agentes; no se imponen de forma exógena.
\end{frame}

\begin{frame}{Características de cada agente}
\textbf{Cada agente $i$ se caracteriza por:}
\begin{itemize}
    \item \textbf{Creencias subjetivas:} $\mu_i \in \mathbb{R}^n$, que representa el retorno esperado de contratar con cada agente de la red (columna $i$ de la matriz $M$).
    \item \textbf{Riesgo percibido:} $\Sigma_i \succ 0$, la matriz de covarianza de los riesgos percibidos.
    \item \textbf{Aversión al riesgo:} $\gamma_i > 0$.
    \item \textbf{Restricciones regulatorias:} $\Psi_i \in \{0,1\}^{|J_i|\times n}$, donde $J_i \subseteq [n]$ es el conjunto de agentes con quienes $i$ \textbf{puede} contratar.
\end{itemize}

\vspace{0.3cm}
\footnotesize{En nuestro caso, $\Psi_i$ codifica la regla: \emph{la MiPyME no puede contratar directamente con CORFO/BID}.}
\end{frame}

\begin{frame}{El problema del agente: maximización sujeta a restricciones}
\textbf{Cada agente $i$ resuelve:}
\begin{equation}
\boxed{\max_{w_i \in \mathbb{R}^n} \; g_i(w_i,P) = w_i^\top(\mu_i - Pe_i) - \gamma_i\, w_i^\top \Sigma_i w_i \quad \text{s.a. } \; \Psi_i w_i = w_i}
\end{equation}

\vspace{0.2cm}
\textbf{Interpretación económica:}
\begin{itemize}
    \item \textbf{Retorno esperado:} $w_i^\top(\mu_i - Pe_i)$, con $w_i$ el vector de portafolio (contratos) del agente.
    \item \textbf{Penalización por riesgo:} $\gamma_i\, w_i^\top \Sigma_i w_i$, donde un $\gamma_i$ mayor implica un agente más averso al riesgo.
    \item \textbf{Restricción:} $\Psi_i w_i = w_i$ fuerza a cero los contratos prohibidos (p.ej. MiPyME$\leftrightarrow$CORFO).
\end{itemize}
\end{frame}

\begin{frame}{Solución óptima y red estable}
\textbf{Portafolio óptimo (Teorema 1, Jalan et al. 2024):}
$$w_i^* = Q_i(\mu_i - Pe_i), \qquad Q_i = \Psi_i^\top\left(2\gamma_i \Psi_i \Sigma_i \Psi_i^\top\right)^{-1}\Psi_i$$

\vspace{0.3cm}
\textbf{Condiciones de estabilidad:}
\begin{itemize}
    \item \textbf{Factibilidad:} $W=W^\top$, $P=-P^\top$, se respetan las restricciones $\Psi_i$.
    \item \textbf{Estabilidad:} cada agente maximiza su utilidad dado un $(W,P)$ factible.
\end{itemize}

\vspace{0.3cm}
La red se alcanza mediante un \textbf{algoritmo de negociación bilateral iterativa} (Algorithm 1) que actualiza los precios $P(t)$ hasta convergencia, con parámetro de amortiguación $\eta$.
\end{frame}

% ===================================================================
% SECCIÓN 3: EXTENSIÓN DUAL (VENEGAS, 2025)
% ===================================================================
\section{Extensión Dual (Venegas, 2025)}

\begin{frame}{La extensión: dos redes en competencia}
Venegas (2025) extiende el modelo agregando una \textbf{segunda red} hacia la que los agentes pueden reasignar su portafolio.

\vspace{0.3cm}
\textbf{Idea central:} el portafolio total de cada agente es \textbf{fijo} y se reparte entre la red Tradicional (T) y la red Alternativa (A):
\begin{equation}
\boxed{w_i^* = w_i^T + w_i^A}
\end{equation}

\vspace{0.2cm}
\begin{itemize}
    \item Fija la exposición total, aislando el efecto de competir por infraestructura.
    \item Refleja restricciones de capital (reasignación, no expansión ilimitada).
    \item Reduce el problema a una única variable de optimización por agente.
\end{itemize}
\end{frame}

\begin{frame}{Maximización sujeta a restricciones: caso dual}
\textbf{Cada agente $i$ resuelve conjuntamente sobre ambas redes:}
\begin{equation}
\boxed{\max_{w_i^T,\, w_i^A} \; g_i^T(w_i^T,P^T) + g_i^A(w_i^A,P^A) \quad \text{s.a. } \; w_i^T + w_i^A = w_i^*,\;\; \Psi_i^T w_i^T = w_i^T,\;\; \Psi_i^A w_i^A = w_i^A}
\end{equation}

\vspace{0.2cm}
\textbf{Utilidad Tradicional:}
$$g_i^T = (w_i^T)^\top(\mu_i - Pe_i) - \gamma_i (w_i^T)^\top\Sigma_i w_i^T$$

\textbf{Utilidad Alternativa} (con beneficio $\tau$ y fricción cuadrática $\lambda_A$):
$$g_i^A = (w_i^A)^\top(\mu_i + \textcolor{SafeGreen}{\tau\mathbf{1}} - Pe_i) - \gamma_i(w_i^A)^\top\Sigma_i w_i^A - \textcolor{AlertRed}{\lambda_A\|w_i^A\|^2}$$
\end{frame}

\begin{frame}{Solución óptima del problema dual}
Resolviendo el problema anterior se obtiene la asignación a la red Tradicional:
\begin{equation}
\boxed{w_i^{T*} = Q_i^{\text{dual}}\cdot a_i}
\end{equation}

\vspace{0.2cm}
\begin{columns}[T]
    \column{0.5\textwidth}
    \textbf{Matriz Q dual:}
    $$Q_i^{\text{dual}} = \left(2\gamma_i\Sigma_i + \textcolor{AlertRed}{\lambda_A I}\right)^{-1}$$
    \footnotesize{(caso simplificado $\Psi^T=I_n$)}

    \column{0.5\textwidth}
    \textbf{Vector de offset:}
    $$a_i = -\tau\mathbf{1} + 2\gamma_i\Sigma_i w_i^* + 2\lambda_A w_i^*$$
\end{columns}

\vspace{0.3cm}
$w_i^{A*} = w_i^* - w_i^{T*}$ se obtiene como residuo. \\
\footnotesize{Validado numéricamente contra el benchmark de Venegas: recuperación de bienestar del 100\% frente al caso sin restricciones.}
\end{frame}

% ===================================================================
% SECCIÓN 4: CASO DE APLICACIÓN AL PROGRAMA BID/CORFO
% ===================================================================
\section{Caso de Aplicación: BID/CORFO}

\begin{frame}{Mapeo institucional: 4 agentes, 2 redes}
\textbf{Estructura del caso real:}
\begin{itemize}
    \item[1] \textbf{CORFO/BID}, que actúa como banca de segundo piso.
    \item[2] \textbf{Intermediario financiero (IFNB)}, ya sea cooperativa, leasing o factoring.
    \item[3] \textbf{MiPyME}, el beneficiario final.
    \item[4] \textbf{Banco tradicional}, el canal de crédito comercial estándar.
\end{itemize}

\vspace{0.3cm}
\begin{block}{Punto clave}
La \textbf{MiPyME (agente 3) es el único agente presente en ambas redes}, y por tanto el único que resuelve el problema dual $w_3^T+w_3^A=w_3^*$. Los agentes 1, 2 y 4 participan en una sola red cada uno.
\end{block}
\end{frame}

\begin{frame}{Red 1 (Tradicional, T): CORFO -- Intermediario -- MiPyME}
\begin{columns}[T]
    \column{0.55\textwidth}
    \textbf{Estructura} (idéntica al Example 5 de Jalan et al.):
    \begin{itemize}
        \item Agentes: CORFO (1), Intermediario (2), MiPyME (3).
        \item Arista \textbf{CORFO -- MiPyME prohibida}: $\Psi_1, \Psi_3$ excluyen ese contrato.
        \item El Intermediario (2) es el único agente sin restricciones ($\Psi_2=I_3$): actúa como \textbf{cuello de botella obligado}.
    \end{itemize}
    \column{0.4\textwidth}
    \centering
    \begin{tikzpicture}[scale=0.8, transform shape]
        \node[circle, draw=DeepBlue, very thick, minimum size=1.3cm, fill=white] (C) at (0,3) {\textbf{CORFO}};
        \node[circle, draw=DeepBlue, very thick, minimum size=1.3cm, fill=white] (M) at (5,3) {\textbf{MiPyME}};
        \node[circle, draw=DeepBlue, very thick, minimum size=1.3cm, fill=gray!20] (I) at (2.5,0.3) {\textbf{Interm.}};
        \draw[AlertRed, line width=2.2pt, dashed] (C) -- (M) node[midway, above, font=\small, text=AlertRed] {Prohibido};
        \draw[->, line width=2pt, DeepBlue] (C) -- (I) node[midway, left, font=\scriptsize] {$W_{12}$};
        \draw[->, line width=2pt, DeepBlue] (I) -- (M) node[midway, right, font=\scriptsize] {$W_{23}$};
    \end{tikzpicture}
\end{columns}
\end{frame}

\begin{frame}{Red 2 (Alternativa, A): Banco -- MiPyME}
\begin{columns}[T]
    \column{0.55\textwidth}
    \textbf{Estructura} (caso simple de 2 firmas, análogo al Example 4 de Jalan et al.):
    \begin{itemize}
        \item Agentes: Banco tradicional (4), MiPyME (3).
        \item \textbf{Un solo par, sin restricciones}: $\Psi_4=\Psi_3^A=I$.
        \item Canal directo: la MiPyME accede a crédito bancario \textbf{sin pasar por un intermediario regulado}.
    \end{itemize}
    \column{0.4\textwidth}
    \centering
    \begin{tikzpicture}[scale=0.8, transform shape]
        \node[circle, draw=DeepBlue, very thick, minimum size=1.3cm, fill=white] (B) at (0,1.5) {\textbf{Banco}};
        \node[circle, draw=DeepBlue, very thick, minimum size=1.3cm, fill=white] (M) at (5,1.5) {\textbf{MiPyME}};
        \draw[->, line width=2pt, SafeGreen] (B) -- (M) node[midway, above, font=\scriptsize] {$W_{4}=W_{43}$};
    \end{tikzpicture}
\end{columns}

\vspace{0.4cm}
\footnotesize{Nota de nomenclatura: en la terminología T/A del modelo, ``Alternativa'' no implica informal ni no regulada: aquí distingue el canal de mercado (Banco) del canal público-subsidiado (CORFO--Intermediario).}
\end{frame}

\begin{frame}{El problema de maximización aplicado al caso}
Sólo la MiPyME (agente 3) resuelve el problema dual completo:
\begin{equation}
\max_{w_3^T,\, w_3^A} \;\; g_3^T(w_3^T, P^T) + g_3^A(w_3^A, P^A) \quad \text{s.a.} \quad w_3^T + w_3^A = w_3^*
\end{equation}
sujeto a $\Psi_3^T$ (excluye a CORFO) y $\Psi_3^A=I$ (sin restricción con el Banco).

\vspace{0.3cm}
Los agentes 1 (CORFO), 2 (Intermediario) y 4 (Banco) resuelven el problema \textbf{estándar de red única} (Teorema 1), cada uno sobre su propia red:
$$w_1^* = Q_1(\mu_1-P^Te_1), \quad w_2^*=Q_2(\mu_2-P^Te_2), \quad w_4^*=Q_4(\mu_4-P^Ae_4)$$

\vspace{0.2cm}
\footnotesize{Consecuencia práctica: la capa de acoplamiento T/A del modelo dual sólo necesita resolverse para un agente, simplificando la calibración numérica.}
\end{frame}

% ===================================================================
% SECCIÓN 5: DATOS Y FUENTES
% ===================================================================
\section{Datos y Fuentes}

\begin{frame}{Qué necesita el modelo}
\begin{table}
\centering
\small
\begin{tabular}{@{}p{2.6cm}p{4.6cm}p{4.4cm}@{}}
\toprule
\textbf{Primitivo} & \textbf{Variable observable} & \textbf{Fuente} \\
\midrule
$\mu_i$ (creencias) & Tasa de interés del programa vs. tasa de mercado (spread) & CORFO (tasa programa) + CMF (tasa mercado) \\
$\Sigma_i$ (riesgo) & Índice de morosidad (NPL) por cartera/región; atrasos en pago & CMF, ELE (módulo Contabilidad y Finanzas) \\
$\gamma_i$ (aversión al riesgo) & Proxy por tamaño/sector de la firma & ELE, literatura de finanzas corporativas \\
$W_{ij}(t)$ (contratos) & Desembolsos CORFO$\to$Interm. ($W_{12}$), Interm.$\to$MiPyME ($W_{23}$), Banco$\to$MiPyME ($W_{4}$) & CORFO/BID; CMF \\
$\Psi_i$ (restricciones) & MiPyME no accede directo a fondos CORFO & Diseño normativo del programa \\
\bottomrule
\end{tabular}
\end{table}
\end{frame}

\begin{frame}{Fuente 1: CORFO / BID (Red Tradicional)}
\begin{itemize}
    \item \textbf{Cuentas Públicas Participativas CORFO} (2014--2024): panel público construido con 6 años de cobertura (2016, 2019, 2020, 2021, 2023, 2024).
    \item Reporta $W_{12}$ (CORFO $\to$ IFNB) y $W_{23}$ (IFNB $\to$ MiPyME) en montos anuales y número de beneficiarios.
    \item \textbf{Solicitud formal a CORFO en curso} (tras reunión con la institución), pidiendo: serie 2014--2024 por IFNB individual, desglosando desembolso nuevo de saldo de cartera vigente, y tasas por tramo.
    \item Antecedente institucional: préstamo BID de US\$120 millones a CORFO (2016) dirigido específicamente a intermediarios financieros no bancarios.
\end{itemize}
\end{frame}

\begin{frame}{Fuente 2: CMF (Red Alternativa)}
\textbf{Comisión para el Mercado Financiero}, fuente de datos del canal bancario tradicional:
\begin{itemize}
    \item \textbf{Colocaciones comerciales}, cartera grupal segmento MiPyme (portal BEST) $\to$ proxy de $W_4$.
    \item \textbf{Tasas de interés por segmento}: diferencial entre tasa de mercado y tasa subsidiada del Crédito Corfo Mipyme $\to$ proxy de $\mu_4$ y del spread $\tau$.
    \item \textbf{Índice de morosidad} por cartera y región, serie mensual desde 2008 $\to$ candidato para estimar $\Sigma$.
\end{itemize}
\end{frame}

\begin{frame}{Fuente 3: Encuesta Longitudinal de Empresas (ELE, INE)}
\begin{itemize}
    \item Encuesta \textbf{bienal} del INE/Ministerio de Economía; última ola publicada: \textbf{ELE-7} (período de referencia 2020--2022).
    \item Su módulo de \textbf{Contabilidad y Finanzas} pregunta, a nivel de empresa individual: préstamos/crédito obtenidos, tasa de interés nominal del crédito de mayor monto, atrasos en pago (mora) y garantías exigidas.
    \item Permite aproximar $\Sigma_3$ y $\mu_3$ (lado MiPyME) sin depender de datos de empresa que CORFO probablemente no compartiría por confidencialidad.
\end{itemize}

\vspace{0.3cm}
\begin{alertblock}{Limitaciones}
No identifica si el crédito proviene específicamente del programa BID/CORFO (sólo categorías generales de fuente); cubre sólo empresas formales con ventas $>$800 UF; periodicidad bienal (menor densidad temporal que la ideal).
\end{alertblock}
\end{frame}

% ===================================================================
% SECCIÓN 6: DIAGNÓSTICO DE IDENTIFICACIÓN
% ===================================================================
\section{Diagnóstico de Identificación}

\begin{frame}{¿Cuántos datos se necesitan? (Teorema 5, Jalan et al.)}
\begin{block}{Teorema 5 --- Tamaños de muestra pequeños bastan}
Si cada agente estima $\hat\Sigma_i$ a partir de $m$ observaciones independientes con $m=\lceil n\log n\rceil$, el umbral de amortiguación estimado $\hat\eta^*$ converge al $\eta^*$ verdadero con alta probabilidad.
\end{block}

\vspace{0.3cm}
Para la Red Tradicional ($n=3$ agentes): $m=4$ (log natural) o $m=5$ (log base 2) observaciones \textbf{conjuntas} por agente.

\vspace{0.2cm}
\footnotesize{Es el resultado del paper más directamente aplicable a la validación empírica: da una cota concreta y defendible de cuántos períodos históricos se necesitan antes de calibrar.}
\end{frame}

\begin{frame}{Diagnóstico del panel actual}
\begin{itemize}
    \item El panel público CORFO cubre \textbf{6 años} (2016--2024), pero sólo \textbf{3} tienen observación conjunta de $W_{12}$ y $W_{23}$ en el mismo año (2019, 2020, 2024).
    \item De esos 3, sólo \textbf{2 años (2019, 2020)} están libres de la mezcla flujo/stock detectada en $W_{23}$ (que incluye cartera rotativa de años previos, no sólo desembolso nuevo).
\end{itemize}

\vspace{0.4cm}
\begin{table}
\centering
\begin{tabular}{lcc}
\toprule
& \textbf{Observado} & \textbf{Mínimo teórico (Teorema 5)} \\
\midrule
Años con $W_{12}$ y $W_{23}$ conjuntos & 3 & \multirow{2}{*}{4--5} \\
Años libres de mezcla flujo/stock & 2 & \\
\bottomrule
\end{tabular}
\end{table}
\centering
\textcolor{OrangeAccent}{\textbf{Cobertura: 50\%--75\% del mínimo teórico}}, lo que motiva la solicitud formal de datos a CORFO.
\end{frame}

\begin{frame}{Limitación estructural: identificabilidad de $\Sigma_{13}$}
\begin{itemize}
    \item La Proposición 1 de Jalan et al. (SDP para inferir $\Sigma$ desde incrementos $\Delta W(t)$) fue diseñada para redes \textbf{sin restricciones} ($\Psi_i=I_n$).
    \item Bajo la arista prohibida CORFO--MiPyME, $W_{13}\equiv 0$ \textbf{siempre}, por lo que nunca aporta información sobre $\Sigma_{13}$.
    \item En un prototipo con datos sintéticos, el error de recuperación de $\Sigma$ no decreció al aumentar el número de períodos observados ($T=15$ a $T=300$).
\end{itemize}

\vspace{0.3cm}
\begin{alertblock}{Hallazgo metodológico propio}
La red regulada, motivación central de la extensión de Venegas, es precisamente el caso en que la Proposición 1 original pierde identificabilidad, lo cual se documenta como limitación en la memoria.
\end{alertblock}
\end{frame}

% ===================================================================
% SECCIÓN 7: ESTADO ACTUAL Y PRÓXIMOS PASOS
% ===================================================================
\section{Estado Actual y Próximos Pasos}

\begin{frame}{Avance a la fecha}
\begin{itemize}
    \item \textbf{Código:} auditoría del modelo de Venegas completada; capa de acoplamiento T/A implementada y validada (62/62 tests) contra el benchmark de recuperación de bienestar del 100\%.
    \item \textbf{Datos:} panel CORFO 2016--2024 construido desde cuentas públicas; solicitud formal de datos desagregados enviada a CORFO.
    \item \textbf{Datos complementarios:} ELE-7 (INE) identificada para calibrar el lado MiPyME sin depender de CORFO.
    \item \textbf{Diagnóstico:} requisito muestral del Teorema 5 cuantificado contra el panel disponible (50--75\% de cobertura).
\end{itemize}
\end{frame}

\begin{frame}{Hoja de ruta}
\begin{block}{Primer hito de avance (27 de septiembre)}
Calibración: construir el panel real BID/CORFO y estimar $M$, $\Sigma_i$ bajo $\Gamma=I$ (aversión al riesgo homogénea), reportando el diagnóstico muestral como hallazgo propio.
\end{block}

\vspace{0.3cm}
\begin{block}{Segundo hito de avance (8 de noviembre)}
Obtención de resultados: correr el modelo dual ya calibrado y obtener contratos, precios y reparto de bienestar, contrastados contra el benchmark de Venegas. La heterogeneidad de agentes ($\gamma_i\neq\gamma_j$) queda como línea tentativa/secundaria.
\end{block}

\vspace{0.3cm}
\footnotesize{Estrategia de contingencia: si la respuesta formal de CORFO se retrasa, se calibra de forma provisional con el panel público + ELE, explícitamente reemplazable cuando lleguen los datos oficiales.}
\end{frame}

\begin{frame}{Conclusiones}
\begin{itemize}
    \item El caso BID/CORFO--Intermediario--MiPyME--Banco mapea de forma natural sobre el modelo dual de Venegas: la MiPyME es el único agente que enfrenta genuinamente la competencia entre infraestructuras de financiamiento.
    \item El modelo se formula como un problema de \textbf{maximización de utilidad media-varianza sujeta a restricciones regulatorias y de portafolio}, resoluble en forma cerrada.
    \item Existen fuentes de datos públicas (CORFO, CMF, ELE) que permiten avanzar en la calibración incluso antes de recibir datos administrativos desagregados.
    \item El propio proceso de calibración revela limitaciones metodológicas del paper original (identificabilidad bajo aristas prohibidas) que constituyen un aporte propio de esta memoria.
\end{itemize}

\vspace{0.5cm}
\centering
\Large{¡Gracias! Preguntas y comentarios son bienvenidos.}
\end{frame}

\end{document}
